% Proposed B.117. Owned fragment; central registers/assembly unchanged.
\subsection*{The actual neck together with its regional base loop}

\statement{B.117}{Proved}{The actual subgroup generated by the common neck and its base loop is an ordinary regional retract with a marked index-two tree graph presentation}
\label{stmt:B117}
Use the actual B.112 regional marking
$\bar x=CxC$, $B=\bar b$, $h=fk$, and $z=Ch$.
Let $R^0$ be the ordinary index-two subgroup of $R$, and let
$K=F(a,B,d,e,f)$ be the actual common neck group from B.114.
Put
\begin{equation}\label{eq:B117-actual-generators}
 L=\langle K,k\rangle,\qquad
 (V_1,V_2,V_3,V_4)=(a,d,e,f^{-1}Bf),\qquad H=\langle h,k\rangle.
\end{equation}
Then $L$ is a retract of $R^0$, with the exact marked presentations
\begin{equation}\label{eq:B117-L-presentation}
 \begin{gathered}
 H=\langle h,k\mid[h^2,k]=1\rangle,\\
 L=H*_{\langle k\rangle}
           \bigl(\langle k\rangle\times F(V_1,V_2,V_3,V_4)\bigr)\\
 \phantom{L}
 =\langle h,k,V_1,V_2,V_3,V_4\mid
       [h^2,k]=1,\ [k,V_i]=1\ (1\leq i\leq4)\rangle.
 \end{gathered}
\end{equation}
In particular $L$ has its full induced profinite topology in $R$.
The actual $K$ is free on $(V_1,V_2,V_3,V_4,f)$, with $f=hk^{-1}$,
and the map
\begin{equation}\label{eq:B117-K-retraction}
 L\longrightarrow K,\qquad k\longmapsto1,\quad h\longmapsto f,
 \quad V_i\longmapsto V_i
\end{equation}
is a retraction.

Define an auxiliary character $\nu:L\to\mathbb Z/2$ by
$\nu(h)=1$ and $\nu(k)=\nu(V_i)=0$, and write $L_2=\ker\nu$.
Then $L_2$ is the graph group on the eleven generators
\begin{equation}\label{eq:B117-tree-generators}
 t=h^2,\quad k_0=k,\quad k_1=hkh^{-1},\quad
 V_{0i}=V_i,\quad V_{1i}=hV_i h^{-1}\quad(1\leq i\leq4),
\end{equation}
with exactly the ten edge relations
\begin{equation}\label{eq:B117-tree-edges}
 [t,k_0]=[t,k_1]=1,\qquad
 [k_0,V_{0i}]=[k_1,V_{1i}]=1\quad(1\leq i\leq4).
\end{equation}
The defining graph is a tree: $t$ is joined to $k_0,k_1$, and
each $k_j$ has its four specified leaves.

The actual subgroup $K_2=K\cap L_2$ is free of rank nine, with
the marked basis
\begin{equation}\label{eq:B117-marked-K2}
 V_{01},\ldots,V_{04},\ V_{11},\ldots,V_{14},\ F=f^2,
 \qquad F=k_1^{-1}t k_0^{-1}.
\end{equation}
Retraction \eqref{eq:B117-K-retraction} restricts to $L_2\to K_2$.
Consequently all the embeddings of $L,L_2,K,K_2$ into their stated
overgroups induce embeddings of their full profinite completions.
The character $\nu$ is auxiliary; no identification with a geometric
double cover is asserted. The word $F$ in
\eqref{eq:B117-marked-K2} is not replaced by the coordinate $t$.
No embedding or centralizer conclusion after B.114's two cross
commutators follows from this theorem.

\begin{proof}
\emph{B.117.1. The full regional graph and its two colors.}
We first give the complete presentation comparison needed for the
retraction, retaining every regional translate. Set
\[
 D=\langle C,k\mid C^2=1\rangle=C_2(C)*\mathbb Z(k),
 \qquad T=\langle k\rangle.
\]
Let $\Omega$ have vertices $(gT,U,n)$ for $gT\in D/T$,
$U\in\{a,b,d,e\}$, and $n\in\mathbb Z$. Its edges are all
$D$-translates, at every $n$, of the ten types
\begin{equation}\label{eq:B117-Omega-edges}
 \begin{gathered}
 (T,U,n)\,--\,(CT,U,n)\quad(U=a,b,d,e),\\
 (T,a,n)\,--\,(T,b,n),\\
 (CT,a,n)\,--\,(T,d,n),\qquad
 (CT,a,n)\,--\,(T,e,n),\\
 (T,d,n)\,--\,(CT,e,n),\\
 (T,b,n)\,--\,(CT,d,n+1),\qquad
 (T,b,n)\,--\,(CT,e,n+1).
 \end{gathered}
\end{equation}
Here $D$ acts by left multiplication on the coset coordinate and
$z$ shifts the integer coordinate. These commuting actions preserve
the graph. Writing $A(\Omega)$ for the group on its vertices with
commutators exactly along its edges, the actual presentation gives
\begin{equation}\label{eq:B117-entire-regional-model}
 R\cong A(\Omega)\rtimes(D\times\mathbb Z_z).
\end{equation}
This is an equality of marked presentations. Indeed, from B.112,
$[C,z]=[k,z]=1$, $C^2=1$, $[k,U]=1$, and the ten node relations
become the edges \eqref{eq:B117-Omega-edges} after putting $h=Cz$.
Map a named vertex to $g z^n U z^{-n}g^{-1}$ in $R$.
Changing the representative of $gT$ does not change this word,
because $k$ commutes with $z^nUz^{-n}$. Every edge maps to an
actual conjugate of a defining node relation. In the reverse direction,
send $C,k,z$ to the acting subgroup and $U$ to $(T,U,0)$.
The stabilizer $T$ and the specified graph edges satisfy all regional
relations. The maps are inverse on every named vertex and on
$C,k,z$. Thus no additional graph relation or old-base translate
has been omitted.

Let $\varepsilon_C:D\to\mathbb Z/2$ send $C$ to $1$ and $k$ to
$0$. Color the vertex $(gT,U,n)$ by
\begin{equation}\label{eq:B117-color}
 \varepsilon_C(g)-n+\mathbf1_{\{U=b\}}\pmod2.
\end{equation}
The representative of $gT$ is immaterial. Each of the four diagonal
types changes coset parity. The $a$--$b$ type changes family offset.
The next three same-level types change coset parity, and each of
the last two types changes coset parity, level parity, and family
offset. Thus every edge joins opposite colors. Translation by $D$
or by a power of $z$ changes both endpoint colors by the same
amount, so the assertion covers the entire graph. Both color classes
induce edgeless full subgraphs.

The original ordinary-cover character kills all graph generators
and takes $(g,z^n)$ to $\varepsilon_C(g)+n$. Its kernel in the
acting group is
\begin{equation}\label{eq:B117-H-kernel}
 H=\{(g,z^n):\varepsilon_C(g)=n\pmod2\}=\langle Cz,k\rangle.
\end{equation}
For the last equality, replace each $C$ in a word for $g$ by $Cz$.
The resulting height has the required parity, and powers of
$(Cz)^2=z^2$ adjust it by any even integer. Since $h=Cz$, this is
the stated $H$. It preserves each color class. Deleting all color-one
graph generators is therefore an $H$-equivariant graph-group
retraction, and extends to
\begin{equation}\label{eq:B117-color-retraction}
 R^0=A(\Omega)\rtimes H
 \longrightarrow A(\Omega_0)\rtimes H.
\end{equation}

\emph{B.117.2. Identify the image with the actual $\langle K,k\rangle$.}
Put $r=Ck^{-1}$, so $f=rz$. The free Schreier basis for the kernel
of the $f$-exponent on the actual $K$ is
$f^n U f^{-n}$ for $U=a,B,d,e$ and $n\in\mathbb Z$.
In \eqref{eq:B117-entire-regional-model} these are literally
\begin{equation}\label{eq:B117-K-vertices}
 \begin{cases}
 (r^nT,U,n),&U=a,d,e,\\
 (r^nCT,b,n),&U=B.
 \end{cases}
\end{equation}
They all have color zero. Also $f\in H$, so the entire $K$ lies
in the image of \eqref{eq:B117-color-retraction}.
The four vertices $a,d,e,f^{-1}Bf$ are
\[
 (T,a,0),\quad(T,d,0),\quad(T,e,0),\quad(T,b,-1),
\]
where the last equality uses $r^{-1}C=k$. Each has stabilizer
exactly $T$ in $H$. Their four $H$-orbits are disjoint and exhaust
$\Omega_0$: for $a,d,e$ the orbit condition is
$\varepsilon_C(g)=n$, and for the $b$ representative at level $-1$
it is $\varepsilon_C(g)=n+1$. These are exactly
\eqref{eq:B117-color} with color zero.

Since the full graph $\Omega_0$ has no edges,
$A(\Omega_0)$ is free on these four $H/T$-orbits. Its semidirect
product with $H$ has the presentation
\[
 \langle H,V_1,V_2,V_3,V_4\mid[k,V_i]=1\ (1\leq i\leq4)\rangle.
\]
Explicit inverse maps send $(gT,i)$ to $gV_i g^{-1}$ in one
direction and $V_i$ to $(T,i)$ in the other, fixing $H$.
The first map is independent of the coset representative precisely
because $[k,V_i]=1$. This proves the amalgam description without
an inference from orbit counts alone.

The subgroup is exactly $L=\langle K,k\rangle$. One inclusion
follows because $K$ contains all four $V_i$ and $h=fk$.
Conversely $a,d,e,f\in\langle H,V_i\rangle$ and
$B=fV_4f^{-1}$, so all five original generators of $K$ lie there.
Consequently \eqref{eq:B117-color-retraction} is the claimed
retraction onto the actual marked $L$.

\emph{B.117.3. The finite presentation of $H$ and its tree subgroup.}
In $D\times\mathbb Z_z$, the subgroup $H$ is generated by
$h=Cz,k$ and satisfies $[h^2,k]=1$.
The abstract group $\langle h,k\mid[h^2,k]=1\rangle$ has an
index-two subgroup, obtained by parity of $h$, with presentation
\[
 \langle t,k_0,k_1\mid[t,k_0]=[t,k_1]=1\rangle,
 \quad t=h^2,\ k_0=k,\ k_1=hkh^{-1}.
\]
Its images are $z^2,k,CkC$ in the actual acting group. The two
words $k,CkC$ freely generate the index-two $C$-parity kernel of
$C_2(C)*\mathbb Z(k)$, while $z^2$ generates an independent
central infinite cyclic group. The presentation map is therefore
injective on this index-two subgroup, and also on the quotient of
order two. This proves the exact presentation of $H$ in
\eqref{eq:B117-L-presentation}.

For $L$, use the same transversal $\{1,h\}$ for $\nu$.
The Schreier generators are exactly
\eqref{eq:B117-tree-generators}. Rewriting $[h^2,k]$ at the two
transversal elements gives $[t,k_0]$ and $[t,k_1]$.
Rewriting each $[k,V_i]$ gives $[k_0,V_{0i}]$ and
$[k_1,V_{1i}]$. These are all rewritten defining relations.
This proves \eqref{eq:B117-tree-edges}, and hence the exact
eleven-vertex, ten-edge tree graph presentation. It is a finite-index
subgroup presentation, not a finite quotient of the infinite graph.

\emph{B.117.4. The actual retained branch word and full topology.}
Replacing $B$ by $V_4=f^{-1}Bf$ is a free basis change of $K$.
Thus $K=F(V_1,V_2,V_3,V_4,f)$, and its restriction of $\nu$ sends
$f$ to $1$ and every $V_i$ to $0$. Its Schreier basis is
$V_i,fV_i f^{-1},f^2$. The existing relations $[k,V_i]=1$ imply
\[
 fV_i f^{-1}=hk^{-1}V_i kh^{-1}=hV_i h^{-1}=V_{1i},
\]
and direct multiplication gives
\[
 f^2=hk^{-1}hk^{-1}
     =(hk^{-1}h^{-1})h^2k^{-1}=k_1^{-1}t k_0^{-1}.
\]
This proves the marked basis \eqref{eq:B117-marked-K2}.
In particular there is no identification of the embedded $K_2$
with the coordinate subgroup obtained by replacing $F$ by $t$.

Killing $k$ and sending $h$ to the actual $f$ respects every
relator of \eqref{eq:B117-L-presentation}. It fixes
$f=hk^{-1}$ and the four $V_i$, hence fixes the actual $K$.
This proves \eqref{eq:B117-K-retraction}. The map respects $\nu$
and restricts to the stated retraction $L_2\to K_2$.

For clarity, all topology assertions follow from these actual
retractions and finite index. Given a finite-index normal subgroup
of a retract, its inverse image under the retraction is a finite-index
subgroup of the overgroup with the prescribed intersection.
Therefore the retract has its full induced profinite topology and
its completion embeds. Apply this first to $L\leq R^0$ and to
$K\leq L$, and then to $K_2\leq L_2$. Passing to the finite-index
subgroups $R^0\leq R$ and $L_2\leq L$, and taking cores when
necessary, preserves this topology. No identification of the
completion of an infinite-rank kernel is used.

These conclusions retain the entire regional presentation, every
height, and the actual common-edge words. They give no automatic
embedding after the additional cross relations and do not replace
the actual mirror substitutions by a coordinate action on this tree.
\end{proof}

\paragraph{Exact rewriting certificate.}
The independent finite Schreier checker verifies all ten
relator/transversal rewrites, every one of the eleven generators,
both inverse free-word comparisons, all four leaf-conjugation
identities, the actual $f^2$ word, and the retraction. The color and
whole-orbit proofs above cover the entire infinite graph, independently
of this finite arithmetic certificate.
\begingroup\raggedright
\path{code/check_neck_base_schreier24.py},
\path{records/check_neck_base_schreier24.out},
\path{records/neck_base_schreier24.json};
\path{records/actual_neck_base_subgroup_candidate24.md},
\path{work_CE/review_neck_base_subgroup24.txt}.
\par\endgroup
